\scanpage{81}
\mspage{II. 29}
% archive sheet number: 80
% lines 1–2
5. \uline{Réduction du} \uline{th. de monodromie} \uline{pour} \deleted{\unsure{$H^1$}} \inserted{$H^1(\bar X,\mathbb Z_\ell)$} \uline{au cas} \uline{d'une courbe projective lisse} \uline{géom. connexe}.
% lines 3–5
5.1. La suite spectrale de descente \deleted{\unsure{étale}} \deleted{\unsure{nous}} pour un morphisme fini (SGA 4 \uline{VIII}) nous ramène au cas $X$ \uline{normal}.
% lines 6–9
5.2. \inserted{On peut déjà supposer $X$ affine, et} Une extension \deleted{radicielle} \inserted{finie} du corps de base nous ramène au cas où de plus $X$ est un ouvert d'un schéma projectif \deleted{normal} \inserted{$X'$} géom.$^{\mathrm t}$ connexe et géom.$^{\mathrm t}$ normal.
% lines 10–13
5.3. Dans ce cas, 3.2 (qui implique que
\[H^1(\bar X,\mathbb Z_\ell)-H^1(\bar C,\mathbb Z_\ell)\]
est injectif) nous ramène au cas \deleted{\unsure{de dim}} $X'$ est \inserted{une \uline{courbe}} \deleted{\unsure{géom. normale connexe}} projective et lisse.
% lines 14–17
[\uline{NB} \deleted{Si $K=k(t_1,\ldots,t_r)$, il est} On choisit arbitrairement un anneau de valuation $V'$ dans $K'=K(t)$, \inserted{qui} domine $V$ \unsure{\}}; alors l'image de $I'-I$ est d'indice fini dans $\unsure{V|I}$].
% lines 18–23
Dans ce cas, \deleted{la suite} on peut encore supposer que \inserted{les pts de} $X'-X$ sont rat / $K$. Alors \inserted{dans} la suite exacte (2.5.1), l'opération de $\bar I$ sur les $H^{\unsure{1|2}}_x(\bar X,\mathbb Z_\ell)\simeq\mathbb Z_\ell(-1)$ est triviale, donc il reste à prouver que l'action sur $H^2(\bar X',\mathbb Z_\ell)$ est quasi-unipotente, ce qui est la réduction annoncée.
